\documentclass{article}

\usepackage{url,algorithm,amsthm,amsmath,amsfonts,verbatim}
\usepackage{stix}
\newtheorem{lem}{Lemma}
\newtheorem{thm}{Theorem}
\usepackage[toc,page]{appendix}
\newcommand{\Dmax}{D^\mathrm{max}}
\newcommand{\Wmin}{W^\mathrm{min}}
% some traditional definitions that can be blamed on craig barratt
\newcommand{\BEAS}{\begin{eqnarray*}}
\newcommand{\EEAS}{\end{eqnarray*}}
\newcommand{\BEQ}{\begin{equation}}
\newcommand{\EEQ}{\end{equation}}
\newcommand{\BIT}{\begin{itemize}}
\newcommand{\EIT}{\end{itemize}}

% text abbrevs
\newcommand{\eg}{{\it e.g.}}
\newcommand{\ie}{{\it i.e.}}
\newcommand{\ones}{\mathbf 1}

% std math stuff
\newcommand{\reals}{{\mbox{\bf R}}}
\newcommand{\symm}{{\mbox{\bf S}}}

% probability stuff
\newcommand{\Expect}{\mathbf{E}}
\newcommand{\prob}{\mathbf{Prob}}
\newcommand{\var}{\mathop{\bf var}}

% convexity & optimization stuff
\newcommand{\Co}{{\mathop {\bf Co}}}
\newcommand{\co}{{\mathop {\bf Co}}}
\newcommand{\Var}{\mathop{\bf var{}}}
\newcommand{\dist}{\mathop{\bf dist{}}}
\newcommand{\argmin}{\mathop{\rm argmin}}
\newcommand{\arginf}{\mathop{\rm arginf}}

\newcommand{\dom}{\mathop{\bf dom}}
\newcommand{\argmax}{\mathop{\rm argmax}}

\newtheorem{theorem}{Theorem}

% alg description --- not much for now!
\newenvironment{algdesc}%
{\begin{quote}}{\end{quote}}

% figures
\def\figbox#1{\framebox[\hsize]{\hfil\parbox{0.9\hsize}{#1}}}

% captions a la sirev
\makeatletter
\long\def\@makecaption#1#2{
   \vskip 9pt
   \begin{small}
   \setbox\@tempboxa\hbox{{\bf #1:} #2}
   \ifdim \wd\@tempboxa > 5.5in
        \begin{center}
        \begin{minipage}[t]{5.5in}
        \addtolength{\baselineskip}{-0.95pt}
        {\bf #1:} #2 \par
        \addtolength{\baselineskip}{0.95pt}
        \end{minipage}
        \end{center}
   \else
        \hbox to\hsize{\hfil\box\@tempboxa\hfil}
   \fi
   \end{small}\par
}
\makeatother

\usepackage[preprint]{neurips_2021}
\usepackage[utf8]{inputenc}
\usepackage[T1]{fontenc}
\usepackage{hyperref}
\usepackage{booktabs}
\usepackage{nicefrac}
\usepackage{microtype}
\usepackage{times}
\usepackage{graphicx}
\usepackage{color}

\graphicspath{{figures/build/}}
\hypersetup{
    pdftitle={Trade Intent Models: A Semantic Interface for Cross-Market Trading and Agentic Strategy Execution},
    pdfauthor={Bill Sun and Alpha.Dev team},
    colorlinks=true,
    urlcolor=blue,
    linkcolor=blue,
    citecolor=blue
}

\title{Trade Intent Models: A Semantic Interface for Cross-Market Trading and Agentic Strategy Execution}
\author{Bill Sun and Alpha.Dev team\\ \url{https://alpha.dev/}}
\date{}

\begin{document}
\maketitle
\small

\begin{abstract}
Trading systems remain fragmented by asset class, venue, execution model, and technical stack. A single economic idea may need to be implemented across broker APIs, centralized exchanges, decentralized protocols, options venues, prediction markets, and on-chain transaction flows, each with its own syntax and operational assumptions. This paper argues that trade intent should be treated as a first-class semantic layer for trading systems. A trade intent model is a domain-specific language that compiles ambiguous human requests into a precise, machine-verifiable, execution-agnostic representation of desired economic action. Like semantic parsing in search or CVXPY in optimization, it separates what the user wants from how the system executes it. The central claim is that a well-designed trade intent model creates value through semantic precision, cross-venue interoperability, strategy composability, and agentic automation, while also serving as the representation layer for an action-oriented workflow in which a user provides only an idea and the platform compiles, validates, and deploys a candidate strategy.
\end{abstract}

\section{Introduction}

The core problem is simple: human trading intent is high-level, but execution infrastructure is fragmented and low-level. Users think in terms of outcomes, exposures, hedges, and constraints. Venues expose order types, product identifiers, chain-specific transaction formats, routing rules, and incompatible APIs. Directly mapping one into the other produces brittle systems.

This gap is becoming more acute because the interface to trading is changing. On the human side, more requests now arrive through chat, voice, and what might be called vibe trading: compressed instructions such as ``get me upside into the close,'' ``hedge this if the market turns,'' or ``buy the best expression of this thesis across venues.'' These are natural forms of human communication, but they are not safe execution formats. They are rich in intent and poor in specification.

On the machine side, a growing number of users are connecting general-purpose or non-specialist agents directly to wallets, broker APIs, and DeFi protocols. These agents can read natural language, call tools, and execute multi-step financial actions, but they often lack a stable semantic layer that separates user meaning from irreversible financial execution. In practice, unstructured prompts become direct action surfaces.

Recent incidents show why this matters. In February 2026, Lobstar Wilde, an autonomous crypto agent built on OpenClaw, misinterpreted a small-transfer request and sent about 52.4 million LOBSTAR tokens to a stranger. In March 2026, the National Internet Finance Association of China warned that OpenClaw-style financial deployments could expose sensitive credentials and trigger erroneous transfers or unintended purchases. These examples point to a common systems problem: the missing layer is often not prediction or contract logic, but the absence of a rigorous intermediate representation between ambiguous language and high-privilege action.

This mismatch creates three persistent failures: natural language is ambiguous, execution interfaces are siloed, and automation without a stable intermediate representation is hard to audit or debug.

The proposed compilation pipeline is:
\begin{center}
\small
user idea or natural-language instruction $\rightarrow$ high-level trade intent $\rightarrow$ \\
execution tasks and policies $\rightarrow$ venue-specific orders and transactions.
\end{center}
This framing shifts the goal from better order entry to a semantic interface for action.

\section{Trade Intent as a Domain-Specific Language}

A trade intent model is a formal intermediate representation of desired economic action. It is best understood as a domain-specific language for trading semantics. Its purpose is not merely to record orders, but to express what the user wants in a form that is precise, canonical, execution-agnostic, and machine-operable.

The key design choice is that the schema should represent economic meaning rather than venue syntax. A useful intent object captures the economic objective, instruments and action types, composition structure, constraints, execution preferences, and provenance across parsing, validation, and lowering. This is what allows similar requests to map to similar structures and allows downstream components to reason about the request algorithmically rather than instruction by instruction.

A natural-language request such as
\begin{quote}
\small
Buy \$10 of NVDAx on Solana, \$5 of PAX Gold on the Ethereum L2 with the best available liquidity between Base and Arbitrum, buy one Yes contract on ``Can NVDA exceed \$200 in February 2026?'', buy one contract on ``2 cuts'' for ``How many times will the Fed cut this year?'', and establish an AAPL June 6 190/200 bull call spread through Alpaca via MCP.
\end{quote}
is economically coherent but operationally heterogeneous. A trade intent model allows the system to treat it as one structured trade vector rather than a bag of unrelated API calls.

\section{Why Trade Intent Models Matter}

Trade intent models create value along four dimensions. First, they improve semantic precision and debuggability by forcing a translation step that can detect underspecification, contradictions, or infeasible requests before any trade is placed. Second, they enable cross-venue interoperability: spot, options, perps, prediction contracts, broker flows, and on-chain actions can be represented through one semantic interface even when the final backends remain heterogeneous. Third, they improve strategy composability: multi-leg packages, routing alternatives, cost-risk tradeoffs, and portfolio-level constraints can be reasoned about jointly. Fourth, they provide the shared object needed for agentic automation, so parsers, validators, simulators, planners, monitors, and agents can interoperate through one contract rather than brittle prompt glue.

Without an intermediate representation, these benefits collapse into black-box automation. With one, failures become localizable: the system can determine whether it misunderstood the language, misconstructed the intent, or mistranslated the final execution instructions.

\section{Hierarchical Execution as an Extension}

Execution itself should be layered. The right analogy is robotics: high-level planning is distinct from low-level motor execution. Trading systems should make the same distinction between high-level trade intent and low-level execution tasks or execution policies.

In the context of this paper, that layered execution stack should be read as an extension of the current repository rather than as a claim about what the present codebase already implements. High-level trade intent is the planning layer. It captures the economically meaningful objective in a way that can be derived from natural language without leaving economically relevant ambiguity, and it provides the right surface for cross-venue reasoning, liquidity aggregation, and policy-aware routing.

A request such as ``buy 100{,}000 shares of AAPL before the close'' or ``buy \$1 million of NVDA call options this week'' should first be represented as a portfolio-level object with explicit notional, time horizon, impact tolerance, urgency, and venue constraints. Only then should it be decomposed into lower-level execution tasks, realized through TWAP, VWAP, broker-native smart routing, on-chain naive slicing, AMM or aggregator execution, or richer low-latency strategies that incorporate market-structure or macro-sensitive signals.

This separation creates a principled place for microstructure-level improvement. The semantic layer is responsible for economic correctness; the execution layer is responsible for realized execution quality. The user specifies the portfolio objective once, and the system remains free to search over better low-level implementations without changing that intent. The current repository does not yet include this planning-and-execution hierarchy; it motivates it.

\section{System View and Repository Boundary}

The long-term objective is not a thin brokerage interface. It is an action-description layer and agent workflow substrate in which a user provides an idea, the system synthesizes a candidate strategy, validates it, and deploys it with one click.

\begin{figure}[t]
    \centering
    \includegraphics[width=0.84\linewidth]{pipeline-figure.pdf}
    \caption{TIM acts as the stable representation layer from idea ingestion to execution-policy selection and venue-native deployment.}
    \label{fig:pipeline}
\end{figure}

At a system level, the pipeline is: (1) idea ingestion; (2) semantic parsing into high-level trade intent; (3) strategy synthesis and task decomposition; (4) validation and policy checks; (5) low-level execution policy selection; and (6) execution compilation and deployment.

The trade intent schema is the stable interface across these stages. The accompanying GitHub repository, \texttt{galpha-ai/intent-kit} (\href{https://github.com/galpha-ai/intent-kit}{github.com/galpha-ai/intent-kit}), should be read as the current reference implementation of the contract layer in this architecture: schema definition, template generation, parsing, validation, and dispatch. That stability also allows parsers, schemas, validators, simulators, adapters, and agents to evolve independently while remaining compatible through the shared representation.

For clarity, the present repository implements only part of the full stack proposed in this paper. It covers the semantic gateway layer: intent schemas, agent-facing templates, parsing, validation, and dispatch. High-level portfolio planning, strategy synthesis, execution-policy search, and low-level algorithmic execution are future layers built on top of that interface.

\section{Design Principles and Open Problems}

A practical trade intent language should satisfy five requirements: semantic completeness, canonical structure, separation of concerns, verifiability, and extensibility. A wrapper translates syntax; a real trade intent model organizes semantics. Several research problems remain open: selecting the right abstraction level for the schema; preserving semantic meaning during lowering; deciding which missing fields can be inferred safely; handling partial fills and cross-venue failure recovery; and constraining agent autonomy through risk, policy, and audit layers. These are not peripheral concerns. They are core design constraints for any deployment-grade trade intent system.

\section{Conclusion}

Trade intent models provide a semantic interface between human trading ideas and heterogeneous execution infrastructure. They transform ambiguous input into a precise, machine-verifiable, and execution-agnostic representation of economic action. Their value is broader than cleaner order syntax. A well-designed trade intent layer enables cross-venue interoperability, makes multi-leg strategies composable, improves debuggability, and supports a principled separation between high-level economic intent and low-level execution policy. In that sense, the trade intent model should be viewed not as a feature, but as infrastructure: the shared semantic substrate that allows fragmented markets, heterogeneous APIs, and agent platforms to interoperate through one language of economic intent.

\end{document}
